% Created 2026-06-04 Thu 21:35
% Intended LaTeX compiler: pdflatex
\documentclass[presentation]{beamer}
\usepackage[utf8]{inputenc}
\usepackage[T1]{fontenc}
\usepackage{graphicx}
\usepackage{longtable}
\usepackage{wrapfig}
\usepackage{rotating}
\usepackage[normalem]{ulem}
\usepackage{amsmath}
\usepackage{amssymb}
\usepackage{capt-of}
\usepackage{hyperref}
\makeatletter\def\input@path{{./}{./shared/}{./shared/themes/}}\makeatother
\usepackage{beamerthemequeer}
\usepackage{shared-preamble}
\usetheme{default}
\author{Matt Alexander}
\date{June 6, 2026}
\title{Hopf-Frobenius Gauge Theories}
\hypersetup{
 pdfauthor={Matt Alexander},
 pdftitle={Hopf-Frobenius Gauge Theories},
 pdfkeywords={},
 pdfsubject={},
 pdfcreator={Emacs 30.1 (Org mode 9.7.11)}, 
 pdflang={English}}
\begin{document}

\maketitle
\begin{frame}[label={sec:orga4b1a99},fragile]{Frobenius Algebras}
 \begin{block}{Definitions}
A \emph{(generalized) Frobenius algebra} \((F, \gen{\cdot, \cdot})\), is a not-necessarily associative nor unital algebra \(F\), with non-degenerate bilinear form \(\gen{\cdot, \cdot}: F \otimes F \to k\), such that
\[
\gen{xy,z} = \gen{x,y z}
\]
for all \(x,y,z \in F\).

\newthought{Classical case: } Restricts to finite-dimensional algebras
\begin{itemize}
\item Can define them in terms of compatible algebraic-coalgebraic structure.
\end{itemize}

\newthought{Integrals: } We think of Frobenius algebras as generalizing algebras of integrable functions
\end{block}
\end{frame}
\begin{frame}[label={sec:org60c3887}]{Frobenius Algebras}
{\renewcommand{\arraystretch}{2.0}
\begin{center}
\begin{tabularx}{\textwidth}{lXXX}
 & \shortstack[l]{Frobenius \\  Algebra} & Pairing & Multiplication\\
\hline
Compact supp. & \(C_c^{\infty}(\mathbb{R}^{n})\) & \(\int fg\) & Pointwise\\
Schwartz func. & \(\mathcal{S}(\mathbb{R}^{n})\) & \(\int fg\) & Pointwise\\
Hilb sp. (+ basis) & \((H, \gen{\cdot, \cdot} \{e_i \})\) & \(\gen{x, y}\) & \(e_i e_j = \delta_{ij} e_i\)\\
Minkowksi sp. (+ basis) & \((\mathbb{R}^{4}, \eta, \{e_i \})\) & \(\eta^{ij} x_{i} x_{j}\) & \(e_i e_j = \eta^{ij} e_{i}\)\\
\end{tabularx}
\end{center}


}
\end{frame}
\begin{frame}[label={sec:orgec43f47}]{Hopf Algebras}
\begin{block}{Definition}
A Hopf algebra is a unital associative \$k\$-algebra, \(H\), with maps 

\medskip

\begin{itemize}
\item \fbox{Comultiplication: } \(\Delta: H \to H \otimes H\)

\item \fbox{Counit: } \(\varepsilon: H \to k\)

\item \fbox{Antipode: } An involution \(S: H \to H\)
\end{itemize}

\bigskip

satisfying the appropriate compatibility conditions between the algebraic and coalgebraic structure.


{\renewcommand{\arraystretch}{1.8}
\begin{center}
\begin{tabularx}{\textwidth}{lXXX}
 & \shortstack{Hopf \\Algebra} & Comult. & Counit\\
\hline
Group algs. & \(kG\) & \(\Delta g = g \otimes g\) & \(\varepsilon(g) = 1\)\\
UE algs. & \(U[\mathfrak{g}]\) & \(\Delta p = 1 \otimes p + p \otimes1\) & \(\varepsilon(p) = 0\)\\
\end{tabularx}
\end{center}
}
\end{block}
\end{frame}
\begin{frame}[label={sec:orgb6b227b}]{Hopf-Frobenius Modules}
A left \emph{Hopf-Frobenius module} is a pair \((H,F)\) where \(H\) is a Hopf algebra, \(F\) is a Frobenius algebra and \(F\) is a left \$H\$-module, such that the following conditions are satisfied:

\bigskip

\begin{itemize}
\item \fbox{Weak Cartan: }  \(\sum \gen{g_{(1)} \cdot x, g_{(2)} \cdot y} = \varepsilon(g) \gen{x,y}\)

\item \fbox{Antipode Condition: }  \(\gen{S(g) \cdot x, y} = \gen{x, g \cdot y}\)
\end{itemize}


{\renewcommand{\arraystretch}{1.5}
\begin{center}
\begin{tabularx}{\textwidth}{lXXX}
 & Hopf & Frobenius & Weak Cartan\\
\hline
Fields & \(U[\operatorname{Der}]\) & Schwartz & \(d(fg) = (df)g + f(dg)\)\\
Geometry & \(k\Lambda\) & Minkowski & \(\eta^{ij}\Lambda^{k}_{i}\Lambda^{\ell}_{j} x_{k }y_{\ell} = \eta^{ij} x_{i} y_{j}\)\\
Gauge & \(kG\) & Appropriate VS & \(\gen{g\cdot x, g\cdot y} = \gen{x, y}\)\\
\end{tabularx}
\end{center}
}
\end{frame}
\begin{frame}[label={sec:org3cea3ee},fragile]{Exterior Derivatives and Connections}
 \begin{block}{Definitions}
Let \(F_{1}, F_{2}\) be Frobenius algebras, \(\{e_i\}_{i \in I}\) be a basis for \(F_2\), and let \(\{ \partial_i : F_{1} \to F_1 \}_{i \in I}\) be a collection of derivations on \(F_1\) indexed by the basis of \(F_2\). 
The  \emph{exterior derivative} of \(f\) with respect to the basis \(\{ e_i \}\) and derivations \(\{\partial_i \}\) is the map 
$$d: F_1 \to F_1 \otimes F_{2}$$
defined by
\[
df = \sum_i \partial_i f \otimes e_i
\]
for all f \(\in\) F\textsubscript{1}.

\newthought{Intuition: } \(F_1\) is a Frobenius algebra of functions, and \(F_2\) is a vector space encoding direction. The algebra \(F_1 \otimes F_2\) plays the role of Kähler differentials \(\Omega^1(F_1)\)
\end{block}
\end{frame}
\begin{frame}[label={sec:orgcd80953}]{Connections and Directional Derivatives}
\begin{block}{Definition}
Let \(F_{1}\) be a fixed Frobenius algebra. 
Given a Frobenius algebra \(F_3\), an \$F\textsubscript{1}\$-\emph{connection on \(F_3\)} is a choice of Frobenius algebra \(F_{2}\) and linear morphism 
\[
\nabla: F_{3} \to F_1 \otimes F_2 \otimes F_3.
\]

\bigskip

When \(F_1\) and \(F_3\) are clear, we will simply refer to \(\nabla\) as the connection.
\end{block}
\begin{block}{Definition}
Let \(\nabla: F_3 \to F_1 \otimes F_2 \otimes F_3\) be a connection, and \(d: F_1 \to F_1 \otimes F_2\) be an exterior derivative. The \emph{covariant derivative with respect to \(\nabla\)} is the map \(D_\nabla: F_1 \otimes F_3 \to F_1 \otimes F_2 \otimes F_3\) defined by
\[
D_{\nabla} (f \otimes\theta)= df \otimes \theta + (f \otimes 1 \otimes 1) \nabla\theta
\]
\end{block}
\end{frame}
\begin{frame}[label={sec:orgcb3c154},fragile]{Physics Examples}
 We consider the Frobenius algebras 

\bigskip

\newthought{Functions: } \(F_f = \mathcal{S}(\mathbb{R}^4)\)

\bigskip

\newthought{Spacetime: } \(F_{ST} = (\mathbb{R}^{4}, \eta)\)

\bigskip

\newthought{Gauge: } \(F_G = (V, \gen{\cdot, \cdot})\)
\end{frame}
\begin{frame}[label={sec:org0f68090},fragile]{Connections on Spacetime}
 A \emph{connection on spacetime} has the form
\begin{align*}
\nabla: F_{ST} \to F_f \otimes F_{ST} \otimes F_{ST}.
\end{align*}
If we take \(F_{ST} = \mathbb{R}^4\),
\[
\nabla e_j = \sum_{ik} \Gamma_{ij}^k(x) \otimes e_i \otimes e_k.
\]
The functions \(\Gamma_{ij}^k(x)\) are the \emph{Christoffel symbols}.


\newthought{Covariant derivative: } 
\begin{align*}
D_\nabla f &= df + \sum_i f_i \otimes \nabla e_i \\
&= \sum_{ij} \partial_j f_i \otimes e_j \otimes e_i + \sum_{ijk} f_i(x) \Gamma_{ji}^k (x) \otimes e_j \otimes e_k.
\end{align*}
\end{frame}
\begin{frame}[label={sec:orgb91f87e},fragile]{Connections on Spacetime}
 Apply the bilinear form to the second tensor factor:

\newthought{Directional derivative: } 
\begin{equation*}
\gen{v,D_\nabla f}_2 = \sum_{ij} v_j \partial_j f_i \otimes e_i + \sum_{ijk} \Gamma^k_{ji} v_j f_i \otimes e_k,
\end{equation*}

\bigskip

The covariant and directional derivatives match those from differential geometry.
\end{frame}
\begin{frame}[label={sec:org7ce0bc4},fragile]{Gauge Connections}
 A \emph{gauge connection} has the form
\begin{align*}
\nabla: F_{G} \to F_f \otimes F_{ST} \otimes F_{G}.
\end{align*}

Let \(\{ e_i \}\) be a basis for \(F_{ST}\) and \(\{ \theta_j \}\) be a basis for \(F_G\).

\newthought{Gauge Potential: } If we expand out in a basis, we can write
\[
\nabla \theta_j = \sum_{ik} A_{ij}^k(x) \otimes e_i \otimes \theta_k.
\]
\(A_{ij}^k(x)\) are the \textbf{gauge field} of the gauge connection

\newthought{Covariant derivative: } 
\begin{align*}
D_\nabla f 
&= \sum_{ij} \partial_j f_i \otimes e_j \otimes \theta_i + \sum_{ijk} f_i(x)  A_{ij}^k(x) \otimes e_j \otimes \theta_k.
\end{align*}
\end{frame}
\begin{frame}[label={sec:orgb2be270},fragile]{Gauge Connections}
 \newthought{Directional derivative: } 
\begin{align*}
\gen{v,D_\nabla f }_2
&= \sum_{ij} v_j \partial_j f_i \otimes \theta_i + \sum_{ijk} v_j f_i(x)  A_{ij}^k(x) \otimes \theta_k.
\end{align*}
\begin{block}{(Hopf-Frobenius) Yang-Mills Theory}
We model these theories  with a collection of Hopf-Frobenius modules

{\renewcommand{\arraystretch}{2.0}
\begin{center}
\begin{tabularx}{\textwidth}{XXX}
 & Frobenius & Hopf\\
Fields & \(\mathcal{S}(\mathbb{C}^{N})\) & \(U[\operatorname{Der}]\)\\
Spacetime & \((\mathbb{R}^{4}, \eta)\) & \(\mathbb{R} \Lambda\)\\
Gauge & \(\mathbb{C}^N\) & \(\mathbb{R}SU(N)\), \(N > 1\) or \(\mathbb{R}U(1)\)\\
\end{tabularx}
\end{center}
}
\end{block}
\end{frame}
\end{document}
